\chapter{集合与常用逻辑用语}
\section{集合}

% ITEM: sets-logic-1-0
\solutionitem{sets-logic-1-0}{例A.1}
写出所有既是整数又满足 \(x^2<5\) 的数。若把“整数”改为“实数”，所得集合还相同吗？
\solutionanswer
整数解为 \(\{-2,-1,0,1,2\}\)。实数解是 \((-\sqrt5,\sqrt5)\)，例如 \(\tfrac12\) 属于后者而不属于前者。

% ITEM: sets-logic-1-1
\solutionitem{sets-logic-1-1}{习题A.1（1）}
设 \(A=\{x\in\mathbb Z:x=2k+1,\ k\in\mathbb Z\}\)，\(B=\{x\in\mathbb Z:x=2k-1,\ k\in\mathbb Z\}\)。证明 \(A=B\)。
\solutionanswer
若 \(x=2k+1\)，则 \(x=2(k+1)-1\)，故 \(A\subseteq B\)。若 \(x=2k-1\)，则 \(x=2(k-1)+1\)，故 \(B\subseteq A\)。两次改写用到的 \(k+1,k-1\) 都是整数。

% ITEM: sets-logic-1-2
\solutionitem{sets-logic-1-2}{习题A.1（2）}
列出 \(\{0,1\}\) 的全部子集。其中哪些以 \(0\) 为元素？
\solutionanswer
全部子集是 \(\varnothing,\{0\},\{1\},\{0,1\}\)。其中 \(\{0\}\) 与 \(\{0,1\}\) 含元素 \(0\)。

\section{集合的运算}

% ITEM: sets-logic-2-0
\solutionitem{sets-logic-2-0}{例A.2}
设 \(A=(-2,1]\)，\(B=[1,3)\)。求 \(A\cap B,\ A\cup B,\ A\setminus B\)，并说明 \(1\) 在三个集合中分别是否出现。
\solutionanswer
交集为 \(\{1\}\)，并集为 \((-2,3)\)，差集为 \((-2,1)\)。两个集合都含 \(1\)，所以交、并都保留 \(1\)，从 \(A\) 中去掉 \(B\) 的元素时则去掉它。

% ITEM: sets-logic-2-1
\solutionitem{sets-logic-2-1}{习题A.2（1）}
设 \(A,B\subseteq\Omega\)。证明 \(\Omega\setminus(A\cap B)=(\Omega\setminus A)\cup(\Omega\setminus B)\)。
\solutionanswer
对每个 \(x\in\Omega\)，它不属于 \(A\cap B\)，恰好是指它不属于 \(A\) 或不属于 \(B\)，也就是它属于右边的并集。因此两边元素相同。

% ITEM: sets-logic-2-2
\solutionitem{sets-logic-2-2}{习题A.2（2）}
设 \(A=\{0,1\}\)，\(B=\{1,2\}\)。列出 \(A\times B\) 的元素，并求 \((A\times B)\cap(B\times A)\)。
\solutionanswer
\(A\times B=\{(0,1),(0,2),(1,1),(1,2)\}\)。要同时属于两个积，两个坐标都必须在 \(A\cap B=\{1\}\) 中，所以交集为 \(\{(1,1)\}\)。

\section{命题与条件}

% ITEM: sets-logic-3-0
\solutionitem{sets-logic-3-0}{例A.3}
设 \(x\) 为实数。“\(x>1\)”是“\(x^2>1\)”的什么条件？增加对 \(x\) 的范围限制，使两者等价。
\solutionanswer
前者是后者的充分条件；\(x=-2\) 说明它不是必要条件。由于 \(x^2>1\) 等价于 \(x>1\) 或 \(x<-1\)，例如限制 \(x\geq0\)，就只剩 \(x>1\) 这一支。

% ITEM: sets-logic-3-1
\solutionitem{sets-logic-3-1}{习题A.3（1）}
设 \(a,b\) 是整数。写出“若 \(a+b\) 是奇数，则 \(a,b\) 中至少一个是奇数”的逆命题与逆否命题，并判断真假。
\solutionanswer
逆命题说：若至少一个是奇数，则和是奇数。取 \(a=b=1\) 即知它为假。逆否命题说：若 \(a,b\) 都是偶数，则和是偶数。写成 \(a=2k,b=2l\)，则和为 \(2(k+l)\)，故逆否命题和原命题都真。

% ITEM: sets-logic-3-2
\solutionitem{sets-logic-3-2}{习题A.3（2）}
对于两个集合 \(A,B\)，证明 \(A\subseteq B\) 的一个充要条件是 \(A\cap B=A\)。
\solutionanswer
若 \(A\subseteq B\)，则 \(A\) 中每个元素都在 \(B\) 中，取交集保留整个 \(A\)。反过来，若 \(A\cap B=A\)，则 \(A\) 中任一元素都是交集的元素，因而属于 \(B\)。

\section{量词}

% ITEM: sets-logic-4-0
\solutionitem{sets-logic-4-0}{例A.4}
比较“每个整数都有比它大的整数”和“有一个整数比所有整数都大”。用量词写出两句话，并分别判断真假。
\solutionanswer
两句依次为 \(\forall n\in\mathbb Z,\ \exists m\in\mathbb Z,\ m>n\) 和 \(\exists m\in\mathbb Z,\ \forall n\in\mathbb Z,\ m>n\)。第一句为真：给定 \(n\) 后取 \(m=n+1\)。第二句为假：任给 \(m\)，只须取 \(n=m\) 就不满足不等式。

% ITEM: sets-logic-4-1
\solutionitem{sets-logic-4-1}{习题A.4（1）}
写出“对每个实数 \(x\)，都有 \(x^2\geq x\)”的否定，并给出一个使否定成立的数。
\solutionanswer
否定是“存在实数 \(x\)，使 \(x^2<x\)”。例如 \(x=\tfrac12\) 时，\(\tfrac14<\tfrac12\)。

% ITEM: sets-logic-4-2
\solutionitem{sets-logic-4-2}{习题A.4（2）}
已知整数集 \(S\) 非空。用量词写出“\(S\) 有最大元素”，再把它的否定译成一句话。
\solutionanswer
原命题为 \(\exists m\in S,\ \forall n\in S,\ n\le m\)。否定为 \(\forall m\in S,\ \exists n\in S,\ n>m\)：任取 \(S\) 中的一个数，在 \(S\) 中都能找到比它更大的数。这里要求 \(m\in S\)，因为讨论的是最大元素。
